\section{Memoryless controllers at arbitrary radius}\label{sec:memoryless-radius}

The one-state case admits an exact controller-by-controller law for every success radius.  Its radius-one specialization recovers the exact one-state threshold.

\begin{theorem}[Memoryless radius law]\label{thm:memoryless-radius}
Let $A$ be any one-state controller and let $\rho\ge0$.  For each
\[
 \mathcal C\in\{\mathcal P,\mathcal J,\mathcal T_{\rm grid},\mathcal G\},
\]
one has
\[
 \boxed{
 \tau_{\mathcal C,\rho}(A)=F_{\mathcal C,\rho}(1)=
 \begin{cases}
 2,&0\le\rho<1,\\[1mm]
 \lfloor\rho\rfloor+3,&\rho\ge1.
 \end{cases}}
\]
In particular, within each of these four classes every memoryless controller has the same least trap size.
\end{theorem}
\begin{proof}
Write $k=\lfloor\rho\rfloor$.  By \eqref{eq:integer-radius}, radius $\rho$ and radius $k$ define exactly the same success event.

First suppose $0\le\rho<1$, so success means coincidence.  On a two-cell edge, agents started at opposite endpoints are forced to exchange endpoints forever and never coincide.  A singleton has only a time-zero success.  Every listed class contains the two-cell path, so the exact value is two.

Now let $k\ge1$.  We first prove a lower bound that does not use memorylessness.  If a connected maze $V$ with starts $x,y$ were a trap with $|V|\le k+2$, failure at time zero and connectedness would give
\[
 k+1\le d_1(x,y)\le d_V(x,y)\le |V|-1\le k+1.
\]
All inequalities are therefore equalities.  A shortest $x$--$y$ path has $|V|-1$ edges and visits every vertex.  Any extra edge joining nonconsecutive vertices of that path would shorten $d_V(x,y)$, so the whole maze graph is that induced path, with $x,y$ as its endpoints.  Compulsory motion forces both first moves inward.  If the path is
\[
 x=v_0,v_1,\ldots,v_k,v_{k+1}=y,
\]
then after one round the agents are at $v_1$ and $v_k$, and
\[
 d_1(v_1,v_k)\le d_V(v_1,v_k)=k-1\le k\le\rho.
\]
Thus every controller succeeds on every connected maze of at most $k+2$ cells.

It remains to give a matching trap for an arbitrary one-state rule.  Write $r(M)\in M$ for the chosen direction at mask $M$ and set
\[
 h=r(\{\E,\W\}),\qquad v=r(\{\N,\Sdir\}).
\]
The unit vectors $h$ and $v$ are perpendicular.  On the corner mask $\{h,v\}$ let
\[
 a=r(\{h,v\}),\qquad \{a,b\}=\{h,v\}.
\]
Then
\[
 r(\{a,b\})=a,\qquad r(\{b,-b\})=b.
\]
Consider
\[
 V_k=\{0,a,b,2b,\ldots,(k+1)b\},
\]
whose path order is $a,0,b,2b,\ldots,(k+1)b$.  This is an induced path with $k+3$ cells.  Start the two agents at $a$ and $kb$.  The endpoint move from $a$ is forced, the corner at $0$ chooses $a$, an interior $b$-arm cell chooses $b$, and the outer endpoint is forced back.  Hence
\[
 (a,kb)\longleftrightarrow(0,(k+1)b).
\]
Both configurations have Manhattan separation $k+1>\rho$, so they form a trap.  The same witness belongs to all four listed maze classes.
\end{proof}

% Abstract a/b preferred directions. Horizontal ellipsis denotes omitted cells.
\begin{figure}[tbp]
\centering
\begin{tikzpicture}[maze,x=9mm,y=9mm]
 \coordinate (z) at (0,0); \coordinate (a) at (0,1);
 \coordinate (b1) at (1,0); \coordinate (b2) at (2,0);
 \coordinate (bk) at (3.6,0); \coordinate (bkp) at (4.6,0);
 \draw[supportA] (a)--(z); \draw[quiet] (z)--(b1)--(b2);
 \draw[quiet,densely dotted] (b2)--(bk); \draw[supportB] (bk)--(bkp);
 \node[cell,label=below left:{$0$}] at (z) {};
 \node[startA,label=left:{$a$}] at (a) {};
 \node[cell,label=below:{$b$}] at (b1) {}; \node[cell,label=below:{$2b$}] at (b2) {};
 \node[startB,label=below:{$kb$}] at (bk) {}; \node[cell,label=below:{$(k+1)b$}] at (bkp) {};
 \node[font=\scriptsize,anchor=west] at (-.25,1.8) {Preferred directions $a\perp b$};
 \node[paneltitle] at (6.1,1.55) {A constant radius gap};
 \node[note] at (6.1,.75) {$(a,kb)\quad\longleftrightarrow\quad(0,(k+1)b)$};
 \node[note] at (6.1,-.05) {$d_1=k+1>\rho$};
\end{tikzpicture}
\caption{Memoryless trap for $k=\lfloor\rho\rfloor\ge1$, with starts at $a$ and $kb$. The axes are preferred directions, not absolute compass labels; the long arm is schematic, with intermediate labels coalescing for $k=1,2$.}
\label{fig:memoryless-l-trap}
\end{figure}


The class restriction in Theorem~\ref{thm:memoryless-radius} concerns only the matching upper witness: the small-maze lower bound holds for every class of connected induced grid mazes, while equality follows for any class containing the controller-adapted path $V_k$.

\begin{corollary}[The original one-state threshold]\label{cor:one-state-radius-one}\label{thm:A}
For the original radius-one model,
\[
 F_{\mathcal P,1}(1)=F_{\mathcal J,1}(1)
 =F_{\mathcal T_{\rm grid},1}(1)=F(1)=4.
\]
\end{corollary}

Thus the value four is the first nontrivial instance of the general law $\lfloor\rho\rfloor+3$, rather than a separate low-state phenomenon requiring its own proof.
